% Protocol section. Included from main.tex before the tracer.
\section{Credit protocol}
\label{sec:protocol}

The protocol decides whether a knowledge message may be credited for a
change in next-step AUC on a given corpus and backbone. It does not
define a predictor. The tracer of Section~\ref{sec:encoding} is the
backbone on which it is run here, and Fig.~\ref{fig:architecture}
summarises the pipeline.

\subsection{Three rungs}
\label{sec:protocol-rungs}
Each contrast is placed on one rung before it is run
(Table~\ref{tab:evidence-types}).

\begin{itemize}
  \item \textbf{(i) Capacity-matched message.} The on and off arms share
  every map; only the message content changes. The question--KC
  incidence twin, which replaces the train-only incidence matrix by
  zeros, is this rung.
  \item \textbf{(ii) Architecture-matched input.} The maps stay and the
  input is replaced. For the gap this means T-zero (every gap set to
  $0$) and T-misaligned (gaps permuted within each learner), with
  T-shuffled (gaps permuted across learners) and T-boundary (only
  $1[\Delta t{=}0]$) as diagnostic variants. This is the
  strongest available control for an input, because the control arm
  keeps a map that receives gradient.
  \item \textbf{(iii) Package.} The component is added together with
  its maps, and the number of added parameters is recorded. A pass here
  credits the package, not its input. The linear
  $\log(1{+}\Delta t)$ branch is this rung: one
  $\mathrm{Linear}(1,128)$ with bias (256 parameters) absent from the
  no-time arm.
\end{itemize}

\begin{table}[t]
\caption{Rungs of the credit ladder and the designated control of each.}
\label{tab:evidence-types}
\centering
\footnotesize
\begin{tabularx}{\textwidth}{@{}>{\raggedright\arraybackslash}p{3.2cm}
>{\raggedright\arraybackslash}p{5.0cm}
>{\raggedright\arraybackslash}X@{}}
\toprule
Rung & Designated control & What it isolates \\
\midrule
(i)~Capacity-matched message & Observed vs.\ zero Q$\leftarrow$KC
incidence, same maps & Incidence content on a fixed fusion \\
(ii)~Architecture-matched input & Real vs.\ zeroed or permuted gaps,
same maps & The aligned gap input, not extra parameters \\
(iii)~Package & Time branch vs.\ no time branch (256 added
parameters) & The timing package as a whole \\
\bottomrule
\end{tabularx}
\end{table}

\subsection{Outcomes and decision rules}
\label{sec:protocol-threshold}
For a fixed set $S$ of paired training seeds, let $\Delta_s$ be the
validation AUC of the on arm minus that of the designated control under
seed $s$. Two rules decide credit at a threshold $\tau$.

\begin{itemize}
  \item \textbf{Seed rule.} Every $\Delta_s\ge\tau$.
  \item \textbf{Interval rule.} The lower end of the two-sided 95\%
  interval of the mean paired difference reaches the threshold,
  $\bar\Delta-t_{0.975,|S|-1}\,s_\Delta/\sqrt{|S|}\ge\tau$, where
  $s_\Delta$ is the sample standard deviation of the paired
  $\Delta_s$.
\end{itemize}

Under either rule a message is \textbf{Credited} if the rule holds,
\textbf{Consistent-positive, not credited} if it fails while every
$\Delta_s>0$, and \textbf{Unsupported} otherwise, including a negative
mean. The middle outcome is part of the protocol: a gain that is
positive on every seed is reported without a separate credit line.

The seed rule is a minimum over seeds, so it becomes stricter as $|S|$
grows; $|S|{=}5$ is therefore part of its definition, and maps are
comparable only at the same $|S|$. The interval rule scales with the
noise of the paired contrast itself, which is the relevant noise for a
seed-paired design \cite{bouthillier2021variance}, and it has no
constant fitted to any corpus. It was fixed after all runs had been
scored. When the two rules disagree, both outcomes are reported.

We use $\tau{=}{+}0.002$ throughout. On XES3G5M learner fold~0 this is
about nine times the seed standard deviation of the two package control
arms ($0.00021$ and $0.00023$), and about half the gap between the
timed tracer and the closest contextual baseline in
Table~\ref{tab:auc-xes3g5m} ($+0.0037$). Recounting
the unrounded per-seed differences (Table~\ref{tab:threshold-sensitivity})
shows that the reference map holds for any seed-rule threshold in
$(+0.00175,+0.00236]$. The value is an operating rule, not a hypothesis
test of $\mu_\Delta{=}0$, and not a registered public standard.

\begin{table}[t]
\caption{Seed-pass counts ($k/5$) if the operating threshold is varied.
Every cell is a recount of printed per-seed $\Delta$val on XES3G5M
fold~0 (Tables~\ref{tab:qkc-zero-time},~\ref{tab:a1-timing},
~\ref{tab:qkc-capacity}). No retraining.
The published qualitative map---package credited; aligned gap not
credited separately; incidence not credited---holds for any threshold
in $(+0.00175,+0.00242]$. At or below $+0.00175$ the aligned-gap
contrast against T-zero would also pass $5/5$; above $+0.00242$ the
timing package itself would not. The $+0.002$ rule remains an operating
convention locked for this ladder, not a public registry.
Applying the same absolute threshold on ASSISTments~2012, Junyi, and
FoundationalASSIST is an unverified assumption: seed-SD was measured on
XES3G5M fold~0.}
\label{tab:threshold-sensitivity}
\centering
\footnotesize
\begin{tabular}{@{}lcccc@{}}
\toprule
Threshold & Package (zero-inc.) & vs.\ T-zero & vs.\ T-misaligned & Incidence \\
\midrule
$+0.0010$ & $5/5$ & $5/5$ & $5/5$ & $0/5$ \\
$+0.0015$ & $5/5$ & $5/5$ & $3/5$ & $0/5$ \\
$+0.0020$ & $5/5$ & $0/5$ & $0/5$ & $0/5$ \\
$+0.0025$ & $2/5$ & $0/5$ & $0/5$ & $0/5$ \\
$+0.0030$ & $0/5$ & $0/5$ & $0/5$ & $0/5$ \\
\bottomrule
\end{tabular}
\end{table}


\subsection{Registering a message}
\label{sec:protocol-register}
Algorithm~\ref{alg:register} is the interface. A new message enters
the ladder only with a declared rung, a designated control, five
paired seeds, and an outcome under both rules. A worked example is in
the repository file \path{docs/REGISTER_MESSAGE.md}.

\begin{algorithm}[t]
\caption{Register a message on the credit ladder.}
\label{alg:register}
\begin{algorithmic}[1]
\STATE \textbf{Input:} message $m$; backbone; corpus; learner fold; seeds $S$, $|S|{=}5$; $\tau{=}{+}0.002$
\STATE Choose rung (i), (ii), or (iii) and the designated control
\STATE If rung (iii), record the number of added parameters
\FOR{each seed $s\in S$}
  \STATE Train the on arm and the control with seed $s$ on the same fold
  \STATE $\Delta_s\gets\mathrm{AUC}_{\mathrm{val}}(\mathrm{on})-\mathrm{AUC}_{\mathrm{val}}(\mathrm{control})$
\ENDFOR
\FOR{each rule $R\in\{$seed, interval$\}$}
  \IF{$R$ holds at $\tau$}
    \STATE outcome$_R\gets$ Credited
  \ELSIF{every $\Delta_s>0$}
    \STATE outcome$_R\gets$ Consistent-positive, not credited
  \ELSE
    \STATE outcome$_R\gets$ Unsupported
  \ENDIF
\ENDFOR
\STATE Report both outcomes with corpus, backbone, and rung; a package credit does not credit the input alone
\end{algorithmic}
\end{algorithm}
